\chapter{État de l'art}
\label{chap:etat}

\section{Le vent dans la couche limite atmosphérique}

\subsection{Profil vertical}
Près du sol, le vent est freiné par la rugosité de la surface. Dans la couche de surface et en
atmosphère neutre, la vitesse croît avec la hauteur selon le profil logarithmique
\begin{equation}
  U(z) = \frac{u_*}{\kappa}\,\ln\!\left(\frac{z}{z_0}\right),
\end{equation}
où $u_*$ est la vitesse de frottement, $\kappa \approx 0{,}4$ la constante de von Kármán et $z_0$
la longueur de rugosité \citep{stull1988}. En pratique on utilise aussi la loi de puissance
$U(z) = U(z_r)\,(z/z_r)^{\alpha}$, dont l'exposant $\alpha$ dépend de la stabilité et du terrain.
Ces lois décrivent mal les situations stables, le relief et les transitions entre la couche limite
et l'atmosphère libre, où le vent tend vers le vent géostrophique imposé par le gradient de
pression à grande échelle.

\subsection{Phénomènes locaux}
Sur le littoral tunisien, la \textbf{brise de mer} --- un vent de mer vers terre dans la journée,
dû au contraste de température terre--mer --- impose un cycle diurne marqué de la direction du vent.
Le relief canalise et accélère l'écoulement, crée des zones abritées et des ascendances. Ces
phénomènes ont des échelles de quelques kilomètres, plus fines que les mailles des modèles globaux.

\subsection{Conservation de la masse}
Pour les mouvements de la couche limite, l'approximation anélastique \citep{ogura1962} remplace
l'équation de continuité complète par
\begin{equation}
  \nabla\cdot(\rho_0\,\vec v) = \frac{\partial(\rho_0 u)}{\partial x}
  + \frac{\partial(\rho_0 v)}{\partial y} + \frac{\partial(\rho_0 w)}{\partial z} = 0,
  \label{eq:anelastique}
\end{equation}
où $\rho_0(z)$ est la densité de référence. Cette relation lie la composante verticale $w$ à la
divergence horizontale, et sert à la fois de contrainte d'apprentissage et de critère (DSO2).

\section{Méthodes classiques de reconstruction}

\subsection{Interpolation spatiale}
La pondération inverse à la distance (IDW) \citep{shepard1968} estime la valeur en un point par
une moyenne des mesures pondérée par l'inverse d'une puissance de la distance. Le krigeage
\citep{cressie1993} choisit les poids qui minimisent la variance d'erreur sous un modèle de
covariance (variogramme) ajusté aux données. Combinées à une loi verticale (profil logarithmique
ou loi de puissance), ces méthodes forment les références auxquelles le modèle est comparé
(chapitre~\ref{chap:resultats}). Elles ignorent la physique de l'écoulement et le relief entre les
points de mesure.

\subsection{Modèles numériques et réanalyses}
Les modèles de méso-échelle comme WRF \citep{skamarock2019} résolvent les équations primitives sur
une grille fine et décrivent correctement brise et relief, mais une simulation coûte des heures de
calcul. Les \textbf{réanalyses} combinent un modèle et toutes les observations disponibles :
ERA5 \citep{hersbach2020} fournit l'état de l'atmosphère heure par heure depuis 1940 sur une
grille de 0,25° (environ 25\,km), et ERA5-Land \citep{munozsabater2021} la surface à 0,1°. ERA5
n'est disponible qu'avec environ cinq jours de retard. En temps réel et en prévision, on dispose
des modèles opérationnels, comme le GFS de la NOAA (0,25°, échéances horaires puis triho\-raires) et
son ensemble GEFS.

\section{Apprentissage profond pour les champs physiques}

\subsection{Réseaux convolutifs et graphes}
Les réseaux convolutifs, en particulier l'U-Net \citep{ronneberger2015}, sont devenus la référence
pour les problèmes de type image-à-image, dont la descente d'échelle des champs météorologiques.
Les réseaux de neurones sur graphes traitent naturellement des capteurs irrégulièrement répartis ;
GraphCast \citep{lam2023} les a portés à la prévision météorologique globale.

\subsection{Opérateurs neuronaux}
Un opérateur neuronal apprend une application entre espaces de fonctions plutôt qu'entre vecteurs
de taille fixe. DeepONet \citep{lu2021} sépare une branche qui encode la fonction d'entrée et un
tronc qui encode la position de sortie. L'\textbf{opérateur neuronal de Fourier} (FNO)
\citep{li2021fno} paramètre le noyau d'intégration directement dans l'espace de Fourier :
\begin{equation}
  (\mathcal{K}v)(x) = \mathcal{F}^{-1}\!\big(R_\phi \cdot \mathcal{F}(v)\big)(x),
\end{equation}
où $R_\phi$ est une matrice de poids complexes appliquée aux $k_{\max}$ premiers modes. Chaque
couche voit tout le domaine (champ récepteur global), le coût est celui d'une FFT, et le modèle
apprend des solutions d'équations aux dérivées partielles avec une précision qui a ouvert la voie à
FourCastNet \citep{pathak2022} en prévision globale.

\subsection{Apprentissage contraint par la physique}
Les réseaux informés par la physique (PINN) \citep{raissi2019} ajoutent à la perte le résidu des
équations du problème. Le PINO \citep{li2024pino} combine opérateur neuronal et perte physique.
Dans notre cas, la physique entre de deux façons : par la perte (résidu de masse discret,
condition au sol) et par la construction (composante $w$ recalculée par continuité dans le modèle
livré, section~\ref{sec:wcont}).

\section{Architectures explorées dans le projet}
La première moitié du projet a exploré six familles d'architectures 2D, chacune reconstruisant le
champ horizontal à hauteur de capteur, l'extrapolation verticale étant faite séparément par des
lois de profil (tableau~\ref{tab:archi}).

\begin{table}[H]
\centering
\caption{Architectures explorées puis abandonnées.}
\label{tab:archi}
\begin{tabularx}{\linewidth}{@{}lX@{}}
\toprule
Architecture & Raison de l'abandon\\
\midrule
FNO 2D (trois variantes) & Champ horizontal seulement ; perte de masse 2D sans le terme $\partial(\rho w)/\partial z$ ni correction du relief\\
U-Net / FNO & Même limite 2D ; pas de gain mesuré sur le FNO 2D\\
GNN & Entraînement non abouti (aucun point de contrôle exploitable)\\
DeepONet & Entraîné, mais sortie 2D et coût d'évaluation point par point\\
ResNet, ViT & Références sans structure physique ; ViT non abouti\\
\bottomrule
\end{tabularx}
\end{table}

Deux constats ont décidé du changement de direction. D'une part, la contrainte de masse calculée
en 2D omettait le terme vertical de l'équation~\eqref{eq:anelastique} et traitait un pas
horizontal sur un terrain en pente comme un pas à altitude constante : le critère DSO2 n'était pas
mesuré sur la bonne équation. D'autre part, le générateur de données a été étendu pour produire des
cibles 3D sur douze niveaux : un opérateur supervisé directement en 3D devenait possible. Nous
avons donc convergé sur une seule architecture, le FNO 3D contraint par la physique, décrit au
chapitre~\ref{chap:modele}.

\section{Positionnement}
PC-FNO3D se situe entre les méthodes d'interpolation, rapides mais aveugles à la physique, et les
modèles numériques, physiques mais coûteux. Il reprend l'efficacité du FNO (une passe en moins
d'une seconde sur GPU), y ajoute la physique de la couche limite par la perte et par la
construction de $w$, et, dans sa version réelle, apprend une \emph{correction} du profil ERA5 plutôt
que le champ entier : la grande échelle vient de la réanalyse (ou de la prévision GFS), le réseau
apporte le relief, le littoral et les mesures locales.
